% !TEX root = template.tex

\section{Processing Pipeline}
\label{sec:processing_architecture}

\fig{fig:pipeline} shows the path from a raw CSI recording to an activity
label, from the Doppler front-end to the fusion of the four antenna scores.
Learning acts on the shaded blocks only, and every configuration reuses the
same front-end, windowing, sampler, encoder, fusion rule and harness, so a
difference between two rows is attributable to the training objective and not
to incidental engineering. \textbf{One configuration is deliberately outside
this rule}: C0, the reimplementation of the reference SHARP design, runs on
its own split protocol, five-class label set, backbone and majority-vote
fusion, because it anchors our harness against the prior approach rather than
entering the objective comparison (\secref{sec:res_baseline}); the
architecture axis is isolated instead by an ablation that changes the backbone
\emph{only}.

\begin{figure}[t]
\centering
\begin{tikzpicture}[
  font=\scriptsize,
  >={Latex[width=1.2mm]},
  box/.style={draw, rounded corners=1.5pt, align=center, inner sep=2pt,
              minimum height=8mm, text width=14.5mm},
  sh/.style={box, fill=black!8},
  node distance=2.1mm]
\node[box] (csi) {CSI\\4 antennas};
\node[box, right=of csi] (fe) {Doppler\\front-end};
\node[box, right=of fe] (map) {$340{\times}100$\\map/antenna};
\node[box, right=of map] (win) {windowing\\$+\;\mu,\sigma$};
\node[box, below=7mm of csi] (aug) {augment.\\(train only)};
\node[sh, right=of aug] (enc) {encoder\\$f_\theta$};
\node[sh, right=of enc] (head) {task head\\$+$ loss};
\node[box, right=of head] (fus) {antenna\\fusion};
\draw[->] (csi) -- (fe);
\draw[->] (fe) -- (map);
\draw[->] (map) -- (win);
\draw[->] (aug) -- (enc);
\draw[->] (enc) -- (head);
\draw[->] (head) -- (fus);
\draw[->] (win.south) -- ++(0,-2.4mm) -- ($(aug.north)+(0,2.4mm)$) -- (aug.north);
\draw[->] (fus.south) -- ++(0,-2.6mm) node[below, font=\scriptsize] {label};
\end{tikzpicture}
\caption{Processing pipeline. Split and $(\mu,\sigma)$ are frozen before
training; the prediction is the argmax of the softmax averaged over the four
antennas.}
\label{fig:pipeline}
\end{figure}

\section{Signals and Features}
\label{sec:model}

\textbf{Measurement setup.} The SHARP micro-Doppler
dataset~\cite{Meneghello-2023} was collected on a commodity IEEE~802.11 link
with \emph{identical hardware throughout}: domains differ by room, monitor
position, subject, day and line-of-sight condition, never by radio chipset. It
holds three environments (bedroom, living room, laboratory), three subjects
and seven \emph{capture sets} --- one (room, monitor position, subject,
propagation) configuration each --- over twelve campaigns; the unit of
recording is a \emph{trace}, acquired simultaneously on four antennas, and our
inventory counts $101$ traces and $408$ per-antenna streams. Activities carry
the dataset's letter codes: the five-class core is E
(empty room), W (walking), R (running), J (jumping) and L (sitting still), and
the extended set adds C (sit-down and stand-up) plus H and S, two activities
whose letter-to-activity assignment the released package does not document ---
we keep its codes rather than attach a gloss we cannot source.

\textbf{Front-end, windowing and normalization.} The data is distributed as
Doppler traces already produced by the reference front-end (CSI phase
sanitization, channel reconstruction, and a short-time Fourier transform along
the packet axis with hop $T_c \approx 6$\,ms), giving per antenna a map of
reflected power over $100$ velocity bins and time; we consume it as-is, since
keeping the front-end fixed is what isolates the variable under study. Each
map is cut into windows of $340$ time steps $\times$ $100$ velocity bins
($\approx2.0$\,s), with training stride $100$ ($\approx71\%$ overlap) and
validation/test stride $340$ so that evaluation units are \emph{disjoint}: the
primary rotation trains on $\approx69$k (window, antenna) samples and is
tested on $425$ fused windows. Two global scalars $(\mu,\sigma)$, computed on
training windows only, are stored in the frozen split file and reused
unchanged elsewhere. Both axes carry a physical orientation, so flipping
velocity (it inverts the Doppler sign) and reversing time (it maps sit-down
onto stand-up) are forbidden augmentations.

\textbf{Augmentation.} \tab{tab:augment} lists the on-the-fly pipeline,
applied after standardization in a fixed order, masks filled with the
post-standardization mean. The cross-entropy configurations use the moderate
column; the contrastive one draws \emph{two independent views} of each window
with the stronger probabilities of the last column, since SupCon needs views
hard enough for its positives to be informative. Velocity is masked sparingly,
as that axis carries the feature separating walking from running.

\begin{table}[t]
\centering
\caption{Augmentation pipeline, in application order after standardization.
The last two columns give the probability for the cross-entropy profile and
for each of the two contrastive views.}
\label{tab:augment}
\resizebox{\columnwidth}{!}{%
\begin{tabular}{lllcc}
\hline
\textbf{Transform} & \textbf{Axis} & \textbf{Parameters} & \textbf{CE} & \textbf{SupCon view} \\
\hline
Circular shift    & time     & $\pm10$ frames                     & 0.5 & 0.5 \\
Masking           & time     & 1--2 masks, width 5--20            & 0.5 & 0.8 \\
Masking           & velocity & 1--2 masks, width 2--10 bins       & 0.5 & 0.8 \\
Amplitude scaling & ---      & $s \sim \mathcal{U}(0.8,1.2)$      & 0.5 & 0.8 \\
Gaussian noise    & ---      & additive, $\sigma = 0.05$          & --- & 0.5 \\
\hline
\end{tabular}}
\end{table}

\textbf{Splits and rotations.} Two rules are enforced by blocking assertions:
\emph{we split by trace, never by window} --- all windows of a trace, on all
four antennas, fall on one side, since otherwise the static channel leaks and
inflates every metric --- and splits are frozen once as JSON under version
control. The primary rotation holds out the laboratory ($81/9/11$ traces), the
hardest choice available, and a second holds out the living room ($80/6/15$)
for replication; a held-out capture set varies room, monitor position, subject
and day at once, so we speak of held-out \emph{sessions} throughout. Because
the class distribution is very uneven across capture sets, every (capture set,
activity) cell with fewer than four traces has one trace pinned into training;
the price is that no C, E or S trace reaches validation on the primary
rotation, making validation macro-F1 a five-class within-run selection signal
--- common-mode across configurations, and with a per-class effect that
\secref{sec:res_err} measures. A third, leave-the-bedroom-out rotation was
formally rejected before any run: the bedroom holds five of the seven capture
sets, so its training side would be two single-session environments totalling
$26$ traces --- confounding ``the model does not generalize'' with ``the model
had no data'' --- with a two-trace validation split. That same skewed
incidence has a structural consequence (\secref{sec:res_diag}) which the
rejected rotation could only have restated.

\section{Learning Framework}
\label{sec:learning_framework}

\fig{fig:learning} shows how the shared encoder is trained under the three
objectives.

\textbf{Backbones.} C1--C4 share a residual convolutional
encoder~\cite{He-2016} adapted to a $340\times100$ input whose axes are not
interchangeable. A stock ResNet-18 stem would crush the velocity axis to a few
bins; instead the stem (conv $3{\times}3$ and max-pool, both stride $(2,1)$)
and stages 3--4 downsample \emph{time} only, while one isotropic stride
$(2,2)$ in stage 2 halves velocity once, leaving a final map of
$\approx11{\times}50$. The four residual stages are halved in width
($32/64/128/256$ channels) and global average pooling yields the feature
vector, $d_{\mathrm{enc}}=256$; at $4.66$ GFLOPs per window and $2.79$\,M
parameters, the encoder was selected on a throughput and memory gate, not on
accuracy. C0 and the backbone ablation use instead a faithful reimplementation
of the SHARP classifier~\cite{Meneghello-2023}: three parallel shallow
branches concatenated, reduced by a $1{\times}1$ convolution, flattened to
$25\,500$ features and mapped to the classes by one dense layer ---
$\approx1$k convolutional against $\approx204$k dense parameters, a
\emph{near-linear wide} model, the opposite design point.

\textbf{Heads and objectives.} All heads are parametrized on
$d_{\mathrm{enc}}$. The activity head is linear, trained with cross-entropy
and label smoothing $0.1$. The adversarial head is an MLP
($d_{\mathrm{enc}} \to d_{\mathrm{enc}}/2 \to$ 6 capture sets, dropout $0.3$)
behind a gradient reversal layer, giving the C2 objective
\begin{equation}
\mathcal{L} = \mathcal{L}_{\mathrm{act}} + \lambda(p)\,\mathcal{L}_{\mathrm{dom}},
\qquad
\lambda(p) = \frac{2}{1+e^{-10p}} - 1 ,
\label{eq:grl}
\end{equation}
with $p = \min(\mathrm{epoch}/20, 1)$, a ramp fixed in advance so the
adversary reaches full strength independently of when early stopping fires.
The contrastive head is an MLP $d_{\mathrm{enc}} \to d_{\mathrm{enc}} \to 128$
with $\ell_2$-normalized output, trained with the supervised contrastive
loss~\cite{Khosla-2020} at temperature $\tau=0.1$, which pulls together all
views sharing a label against every other view in the batch. C3 trains in two
phases: phase A optimizes it for a fixed horizon of $60$ epochs with no
in-loop selection, checkpointing on a pre-committed grid (epochs $40/50/60$);
phase B freezes the encoder, discards the projection head and fits a linear
probe on cached features, and only the grid checkpoint with the best
validation macro-F1 reaches test. That frozen-probe recipe (Adam, $10^{-3}$,
$\leq30$ epochs, early stopping on the shared metric) is reused for every
probe in this report, so encoders are always compared under one readout.

\begin{figure}[t]
\centering
\begin{tikzpicture}[
  font=\scriptsize,
  >={Latex[width=1.2mm]},
  box/.style={draw, rounded corners=1.5pt, align=center, inner sep=2.5pt,
              minimum height=7mm, text width=17mm},
  enc/.style={box, fill=black!8, text width=19mm, minimum height=13mm},
  tag/.style={font=\scriptsize\itshape}]
\node[box] (in) {window\\$340{\times}100$};
\node[enc, right=4mm of in] (e) {encoder $f_\theta$\\ResNet V-B\\$d_{\mathrm{enc}}{=}256$};
\node[box, right=7mm of e, yshift=10mm] (ce) {linear head\\CE $+$ smoothing};
\node[box, right=7mm of e] (grl) {GRL $\lambda(p)$\\capture-set MLP};
\node[box, right=7mm of e, yshift=-10mm] (sc) {projection head\\SupCon, $\tau{=}0.1$};
\node[tag, right=1.5mm of ce.east]  {C1, C2};
\node[tag, right=1.5mm of grl.east] {C2};
\node[tag, right=1.5mm of sc.east]  {C3};
\draw[->] (in) -- (e);
\draw[->] (e.east) -- ++(3.5mm,0) |- (ce.west);
\draw[->] (e.east) -- (grl.west);
\draw[->] (e.east) -- ++(3.5mm,0) |- (sc.west);
\node[box, below=8mm of e, text width=42mm, xshift=6mm] (pb)
  {phase B: freeze $f_\theta$, cache features, fit linear probe};
\draw[->, dashed] (sc.south) -- ++(0,-3mm) -| (pb.east);
\draw[->, dashed] (e.south) -- (pb.north);
\end{tikzpicture}
\caption{Learning framework: one encoder, three objectives --- cross-entropy
(C1), cross-entropy plus a capture-set adversary behind a gradient reversal
layer (C2), and supervised contrastive pre-training read out by a
frozen-encoder linear probe (C3). The dashed path doubles as the diagnostic
instrument: any encoder can be frozen, cached and probed for the activity or
for a domain attribute.}
\label{fig:learning}
\end{figure}

\textbf{Configurations.} The five configurations span a $2{\times}2$ design
over loss family and adversary, plus the anchor: C0 (prior approach), C1 (CE,
our model and internal control), C2 (CE+GRL), C3 (SupCon) and C4 (SupCon+GRL).
C4 was \emph{closed on evidence and never trained}: the diagnostics of
\secref{sec:res_diag} show no readable domain target under either loss family,
so its expected outcome was ``C3 plus noise''. Three pre-registered arms
complete the picture --- C3-ft, a full CE fine-tune from the SupCon encoder;
C1-aug, the amplitude-augmentation arm; and C1-sharplike, the reference
network inside the exact C1 recipe, only the backbone differing.

\textbf{Batching, optimization and reproducibility.} The CE configurations
draw uniform batches of $256$ (window, antenna) samples and optimize an
\emph{unweighted} loss while being selected on macro-F1, a declared asymmetry
kept for comparability with the reference recipe. The contrastive loss is
intra-batch and cannot use uniform batches: a $P{\times}K$ sampler ($P=8$
classes, $K=32$ traces each) draws each class's traces round-robin over the
capture sets where the class exists, making positives ``same activity,
different domain'' by construction, and admits at most one window per trace
per batch (falling back to a $\geq340$-frame offset for classes with fewer
than $K$ traces), since overlapping windows of one recording would reintroduce
the trivial positives that rule exists to remove. Optimization is AdamW
(weight decay $10^{-4}$, peak learning rate $10^{-3}$, cosine decay, $5$
warm-up epochs), mixed precision, gradient clipping at $1.0$, epochs of $400$
fixed steps ($300$ in phase A), at most $40$ epochs, early stopping on fused
validation macro-F1 with patience $10$. Seed $42$ governs initialization,
sampler reseeding and augmentation; two pre-registered seed-$43$ replicates
measure the training-variance floor.

\textbf{Evaluation harness and the single-use test rule.} One shared harness
maps a checkpoint to per-set metric files --- accuracy, macro-F1 (the mean
over classes of the harmonic mean of precision and recall), per-class F1,
confusion matrices and expected calibration error --- always reported with the
number of test traces, since that is the effective sample size. Checkpoints
are selected on validation only and the \emph{test set is evaluated exactly
once}: one session over a frozen, pre-registered row list logged before any
test data is read, so evaluate-then-decide selection is impossible.
